% Source: https://artofproblemsolving.com/community/c3h1345946_prove_that
% Stable identity: aops-topic-1345946-q01
The term $5n^3(n-1)$ is always a multiple of $5$, so the expression is divisible by $5$ exactly when $n^2+1$ is. The claim fails for $n=1$: the expression is $2+0=2$.

Squares modulo $5$:
\[
n\equiv0,1,2,3,4\pmod5\quad\Longrightarrow\quad n^2\equiv0,1,4,4,1\pmod5 .
\]
So $n^2+1\equiv0\pmod5$ exactly when $n^2\equiv4$, that is when $\boxed{n\equiv2\text{ or }3\pmod5}$. For example $n=2$ gives $5+40=45$, divisible by $5$; $n=1$ and $n=5$ give $2$ and $526$, which are not.
